% Section 6.3, from lectures/L06/README.md: the objectives, the evaluation
% questions, and the next lecture.

\section{Review}
\label{c6:sec:review}

\needspace{6\baselineskip}
\subsection*{What you should be able to do now}
After this chapter, you should:

\begin{itemize}
  \item Be able to add shell commands to an application, and use the shell to inspect a running system.
  \item Be able to store and load settings, and understand where they are stored on \code{native_sim}.
  \item Be able to write TCP and UDP servers with Zephyr’s socket API, and reach them from the host.
  \item Understand how Zephyr handles a fatal error, and be able to decide what happens after one.
  \item Be able to find memory errors and undefined behaviour with the sanitizers.
  \item Be able to debug a \code{native_sim} program with GDB, and measure test coverage.
\end{itemize}

\needspace{6\baselineskip}
\subsection*{Questions to test yourself}
You should be able to explain:

\begin{enumerate}
  \item What does a shell command handler receive, and how does it report an error?
  \item Why can \code{log enable dbg} not show debug messages from a module compiled at level \code{INF}?
  \item Where does the settings subsystem store its data on \code{native_sim}, and what makes a run a first
    boot?
  \item What does the offloaded sockets driver do, and why does it need no network setup or root?
  \item What is the difference between \code{k_oops()} and \code{k_panic()}, and what does
    \code{k_sys_fatal_error_handler()} decide?
  \item Why does AddressSanitizer work on \code{native_sim}, and why is it not available when the same code
    is built for a board?
  \item Why does a UBSan report not stop the program, while an ASan report does?
  \item What would change, and what would not, if the driver from Chapter~\ref{c5:ch} and these subsystems
    moved to a real board?
\end{enumerate}

\needspace{6\baselineskip}
\subsection*{Where to go from here}
This is the last chapter: you have built, configured, described, scheduled, driven, tested and
debugged a Zephyr application, and every one of those programs builds for a real board as it stands.
What is left is the board itself.

\textbf{Try it on hardware.} Any board Zephyr supports will do; two cheap ones its tree already
knows are an STM32 Nucleo (\code{-b nucleo_f401re}) and an ESP32 DevKitC
(\code{-b esp32_devkitc/esp32/procpu}). Three things change in the workspace. You need a
cross-compiler, which \code{west sdk install} fetches, where the host's own GCC has stood in until
now. You need the vendor's HAL: \code{native_sim} needed no Zephyr modules, so the course's
\file{west.yml} imports none, and a real SoC needs its own, such as \code{hal_stm32} or
\code{hal_espressif}. And you flash rather than run: \code{west flash} and \code{west debug}
replace \code{west build -t run}, over the Nucleo's on-board ST-LINK or the DevKitC's USB port.

\textbf{What changes on a board} is the part the simulator was never pretending to be: pins and
clocks have to be configured, the console is a real UART, logging is deferred rather than immediate,
interrupts arrive while your code is running, stack overflows can be caught in hardware, and timing
is real, so a busy loop merely wastes power instead of stopping the world.

\textbf{What does not change} is nearly everything this book is about: the kernel API, devicetree
and Kconfig, the driver model, and your tests, which Twister runs on the board itself with
\code{--device-testing}. That is the point of having written them against an emulator.
